%!TEX root = ../main.tex
\section{Proof of Theorem \ref{thm:lqg-variance-one-step}}
\label{sec:app-proof-lqg-variance-one-step}

We begin with a lemma that computes expectations of certain polynomial functions of Gaussian random variables.
While such identities are broadly covered by Isserlis' theorem (Wick's formula), those results mainly provide
expressions for generic even-order moments. In our analysis, we require more structured, problem-specific forms.
Magnus~\cite{magnus78report-gaussianmoments} derives formulas for expectations of products of quadratic forms of the type
$\mathbb{E}\!\left[\prod_{i=1}^{N}\big(u^\top A_i u\big)\right]$.
Our setting involves a non-homogeneous polynomial, hence we state the following lemma.

\begin{lemma}\label{thm:wz_std_lemma}
For $\xi\sim \mathcal N(0,I_m)$, define
\[
w(\xi):=\sum_{i=1}^m (\xi_i^2-1)^2.
\]
Then for any $u\in\mathbb R^m$ and any $M\in \sym{m}$,
\begin{align}
\E[w(\xi)] &= 2m, \label{eq:w0}\\
\E\big[w(\xi)\,(\xi^\top u)^2\big] &= (2m+8)\,\|u\|^2, \label{eq:w_u}\\
\E\big[w(\xi)\,(\xi^\top M \xi)\big] &= (2m+8)\,\tr M, \label{eq:w_M}\\
\E\big[w(\xi)\,(\xi^\top M \xi)^2\big]
&= (4m + 32)\,\tr(M^2) + (2m + 16)\,(\tr M)^2 + 24 \sum_{i=1}^m M_{ii}^2.
\label{eq:w_MM}
\end{align}
\end{lemma}

\begin{proof}
To prove \eqref{eq:w0}, use $\E[\xi^8]=105,\,\E[\xi^6]=15,\,\E[\xi^4]=3,\,\E[\xi^2]=1$:
\[
\E[w(\xi)] = \sum_{i=1}^m \E[(\xi_i^2-1)^2] = \sum_{i=1}^m (3 - 2 + 1) = 2m.
\]

For \eqref{eq:w_u}, expand and note that cross terms with $u_i u_j$ $(i\neq j)$ vanish by odd-moment cancellation,
so only $u_i^2$ terms remain:
\[
\E\big[w(\xi)(\xi^\top u)^2\big]
= \sum_{i=1}^m \sum_{j=1}^m \E\big[(\xi_i^2-1)^2 \xi_j^2\big]\, u_j^2.
\]
Compute
\[
\E\big[(\xi_1^2-1)^2 \xi_1^2\big]=\E[\xi_1^6]-2\E[\xi_1^4]+\E[\xi_1^2]=15-6+1=10,
\qquad
\E\big[(\xi_i^2-1)^2 \xi_1^2\big]=2\ (i\neq 1),
\]
hence $\E[w(\xi)(\xi^\top u)^2]=\|u\|^2(10+2(m-1))=(2m+8)\|u\|^2$.

For \eqref{eq:w_M}, similarly only diagonal entries contribute:
\[
\E[w(\xi)(\xi^\top M \xi)]
= \sum_{i=1}^m\sum_{j=1}^m \E[(\xi_i^2-1)^2\xi_j^2]\, M_{jj}
= (2(m-1)+10)\tr M
= (2m+8)\tr M.
\]

For \eqref{eq:w_MM}, expand the quadratic form as
\[
(z^\top M z)^2=\sum_{i,j,k,l=1}^m M_{ij}M_{kl}\,z_i z_j z_k z_l.
\]
After cancelling the odd terms, only two index patterns contribute: (i) $i=j$ and $k=l$ (diagonal--diagonal terms), and (ii) $\{i,j\}=\{k,l\}$ with $i\neq j$ (off-diagonal pairings), yielding
\[
(z^\top M z)^2
=\sum_{i,k} M_{ii}M_{kk}\,z_i^2 z_k^2
\;+\;
2\sum_{i<j} M_{ij}^2\,z_i^2 z_j^2
\;+\;(\text{terms with vanishing expectation}).
\]
We then evaluate $\E[w(z)z_i^2 z_k^2]$ by splitting into the cases $i=k$ and $i\neq k$ and using the one-dimensional moments
$\E[z^2]=1,\E[z^4]=3,\E[z^6]=15,\E[z^8]=105$, which gives
\[
\E[w(z)\,z_i^4]=\E[(z_i^2-1)^2 z_i^4]+(m-1)\E[(z_1^2-1)^2]\E[z_i^4]=78+6(m-1),
\]
\[
\E[w(z)\,z_i^2 z_k^2]=\E[(z_i^2-1)^2 z_i^2]\E[z_k^2]+\E[(z_k^2-1)^2 z_k^2]\E[z_i^2]+(m-2)\E[(z_1^2-1)^2]\E[z_i^2]\E[z_k^2]
=20+2(m-2),
\quad (i\neq k),
\]
and similarly
\[
\E[w(z)\,z_i^2 z_j^2]=2m+16\qquad (i\neq j).
\]
Substituting these values into the diagonal--diagonal and off-diagonal contributions and collecting terms yields
\[
\E \big[w(z)\,(z^\top M z)^2\big]
= (4m + 32)\,\tr(M^2) + (2m + 16)\,(\tr M)^2 + 24\sum_{i=1}^m M_{ii}^2,
\]
concluding the proof.
\end{proof}

Next we prove Theorem~\ref{thm:lqg-variance-one-step}.

\begin{proof}
\textbf{Proof of $K$ part.}
Using
\[
s_1 = F s_0 + B\varepsilon_0,\quad
a_0 = K s_0 + \varepsilon_0,\quad
r_0 =-\big(s_1^\top Q_s s_1+a_0^\top Q_a a_0\big),
\]
we can expand $r_0$ as
\[
r_0 = -(x + 2y + z),
\]
where
\begin{align}\label{eq:lqr-decompose-r0}
x := s_0^\top M_{ss} s_0,\quad
y := s_0^\top M_{se}\varepsilon_0,\quad
z := \varepsilon_0^\top M_{ee}\varepsilon_0,
\end{align}
and
\[
M_{ss}:=F^\top Q_sF+K^\top Q_a K,\quad
M_{se}:=F^\top Q_sB+K^\top Q_a,\quad
M_{ee}:=B^\top Q_s B+Q_a.
\]
Note $M_{ss}, M_{ee}$ are symmetric and positive semidefinite. The REINFORCE estimator is
\[
\widehat G_K = \nabla_K\log\pi(a_0\mid s_0)\,r_0
= \Sigma^{-1}\varepsilon_0 s_0^\top\, (-(x+2y+z)).
\]
Since $x$ depends only on $s_0$ and $z$ depends only on $\varepsilon_0$, and $s_0\perp\!\!\!\perp \varepsilon_0$ with
$\E[\varepsilon_0]=0$, only the $y$ term remains in expectation:
\[
\E[\widehat G_K]
=
\E\!\left[\Sigma^{-1}\varepsilon_0 s_0^\top (-2y)\right]
= -2\,\Sigma^{-1}\E[\varepsilon_0\varepsilon_0^\top]\,M_{se}^\top \E[s_0 s_0^\top]
= -2\,M_{se}^\top \Sigma_0
= -2\,(Q_aK+B^\top Q_sF)\Sigma_0.
\]

For the second moment, using $\|AB\|_{\mathrm{F}}^2=\tr(B^\top A^\top A B)$,
\[
\|\widehat G_K\|_{\mathrm{F}}^2
= (\varepsilon_0^\top \Sigma^{-2} \varepsilon_0)\,(s_0^\top s_0)\,(x+2y+z)^2.
\]
By independence and odd-moment cancellation,
\[
\E\big[\|\widehat G_K\|_{\mathrm{F}}^2\big]
=
\E\Big[(\varepsilon_0^\top \Sigma^{-2} \varepsilon_0)\,(s_0^\top s_0)\,(x^2 + z^2 + 2xz + 4y^2)\Big]
=
E_{1,K}+E_{2,K}+E_{3,K}+E_{4,K},
\]
with
\begin{align*}
E_{1,K} &:= \E[\varepsilon_0^\top \Sigma^{-2} \varepsilon_0]\;\E\big[(s_0^\top s_0)\,x^2\big],\\
E_{2,K} &:= \E[s_0^\top s_0]\;\E\big[(\varepsilon_0^\top \Sigma^{-2} \varepsilon_0)\,z^2\big],\\
E_{3,K} &:= 2\,\E\big[(s_0^\top s_0)\,x\big]\;\E\big[(\varepsilon_0^\top \Sigma^{-2} \varepsilon_0)\,z\big],\\
E_{4,K} &:= 4\,\E\!\left[(s_0^\top s_0)\,\E\!\left[(\varepsilon_0^\top \Sigma^{-2} \varepsilon_0)\,y^2\,\middle|\,s_0\right]\right].
\end{align*}
We rewrite each term using Lemma~\ref{lem:IS_shorthand}. First,
\[
\E[\varepsilon_0^\top \Sigma^{-2}\varepsilon_0]
=\mathrm{IS}_{\Sigma}\big(\Sigma^{-2}\big),
\qquad
\E[s_0^\top s_0]=\mathrm{IS}_{\Sigma_0}(I_n).
\]
Also,
\[
\E[(s_0^\top s_0)x^2]
=
\mathrm{IS}_{\Sigma_0}\!\big(I_n,M_{ss},M_{ss}\big),
\qquad
\E[(s_0^\top s_0)x]
=
\mathrm{IS}_{\Sigma_0}\!\big(I_n,M_{ss}\big),
\]
and since $z=\varepsilon_0^\top M_{ee}\varepsilon_0$,
\[
\E\big[(\varepsilon_0^\top \Sigma^{-2}\varepsilon_0)\,z\big]
=
\mathrm{IS}_{\Sigma}\big(\Sigma^{-2},M_{ee}\big),
\qquad
\E\big[(\varepsilon_0^\top \Sigma^{-2}\varepsilon_0)\,z^2\big]
=
\mathrm{IS}_{\Sigma}\big(\Sigma^{-2},M_{ee},M_{ee}\big).
\]

It remains to compute $E_{4,K}$. Since $y=s_0^\top M_{se}\varepsilon_0$,
\[
y^2=\varepsilon_0^\top X(s_0)\varepsilon_0,
\qquad
X(s_0):=M_{se}^\top s_0 s_0^\top M_{se}.
\]
Conditioning on $s_0$ and applying Lemma~\ref{lem:IS_shorthand} with $k=2$,
\[
\E_{\varepsilon_0}\!\left[
(\varepsilon_0^\top \Sigma^{-2}\varepsilon_0)\,y^2
\;\middle|\;s_0\right]
=
\mathrm{IS}_{\Sigma}\big(\Sigma^{-2},X(s_0)\big)
=
\tr(\Sigma\Sigma^{-2})\tr(\Sigma X(s_0))
+2\,\tr(\Sigma\Sigma^{-2}\Sigma X(s_0)).
\]
Using $\tr(\Sigma\Sigma^{-2})=\mathrm{IS}_{\Sigma}(\Sigma^{-2})$ and $\Sigma\Sigma^{-2}\Sigma=I$, and defining
\[
U := M_{se}M_{se}^\top,\qquad
V := M_{se}\Sigma M_{se}^\top,
\]
we have
\[
\tr\big(X(s_0)\big) = s_0^\top U s_0,
\qquad
\tr\big(\Sigma X(s_0)\big) = s_0^\top V s_0.
\]
Therefore,
\[
\E_{\varepsilon_0}\!\left[
(\varepsilon_0^\top \Sigma^{-2}\varepsilon_0)\,y^2
\;\middle|\;s_0\right]
=
\mathrm{IS}_{\Sigma}\big(\Sigma^{-2}\big)\,(s_0^\top V s_0)
+2\,(s_0^\top U s_0).
\]
Multiplying by $(s_0^\top s_0)=(s_0^\top I_n s_0)$ and taking expectation over $s_0$ gives
\[
E_{4,K}
=
4\Big(
\mathrm{IS}_{\Sigma}\big(\Sigma^{-2}\big)\;
\mathrm{IS}_{\Sigma_0}\big(I_n,V\big)
+
2\,\mathrm{IS}_{\Sigma_0}\big(I_n,U\big)
\Big).
\]

\medskip
\textbf{Proof of log-std part.}
Write $\varepsilon_0=\Sigma^{1/2}\xi$ with
$\xi\sim\mathcal N(0,I_m)$ and define
\[
q := \Sigma^{-1}(\varepsilon_0\odot \varepsilon_0) - \mathbf 1
= \xi\odot\xi-\mathbf 1,
\qquad
S:=\Sigma^{1/2}M_{ee}\Sigma^{1/2}.
\]
Since $s_0$ and $\varepsilon_0$ are independent,
\[
\E[\widehat G_\ell]
= \E[q\,r_0]
= -\,\E[q\,z],
\qquad
z=\varepsilon_0^\top M_{ee}\varepsilon_0=\xi^\top S\xi.
\]
For the $i$-th component,
\[
\big[\E(q\,z)\big]_i
= \sum_{j,k} S_{jk}\,\E\big[(\xi_i^2-1)\,\xi_j\xi_k\big].
\]
By independence,
\[
\E\big[(\xi_i^2-1)\,\xi_j\xi_k\big]
=\begin{cases}
\E[\xi_i^4-\xi_i^2]=3-1=2, & j=k=i,\\
0, & \text{otherwise},
\end{cases}
\]
so $\big[\E(q\,z)\big]_i = 2\,S_{ii}$ and hence
\[
\E[\widehat G_\ell]
= -2\,\operatorname{diag}(S)
= -2\,\operatorname{diag}\!\big(\Sigma^{1/2}M_{ee}\Sigma^{1/2}\big)
= -2\,\operatorname{diag}\!\big(\Sigma M_{ee}\big),
\]
where the last equality uses that $\Sigma$ is diagonal. 

For the second moment, $\widehat G_\ell = q\,r_0$ and $r_0^2=(x+2y+z)^2$.
Using independence and odd-moment cancellation,
\begin{align*}
\E\big[\|\widehat G_\ell\|_2^2\big]
&= \E\big[\|q\|_2^2\big]\;\E[x^2]
+ \E\big[\|q\|_2^2 z^2\big]
+ 2\,\E[x]\;\E\big[\|q\|_2^2 z\big]
+ 4\,\E\big[\|q\|_2^2 y^2\big].
\end{align*}
The $x$-moments follow from Lemma~\ref{lem:IS_shorthand}:
\[
x=s_0^\top M_{ss}\,s_0,
\qquad
\E[x]=\mathrm{IS}_{\Sigma_0}\!\big(M_{ss}\big),
\qquad
\E[x^2]=\mathrm{IS}_{\Sigma_0}\big(M_{ss},M_{ss}\big).
\]
For the $\varepsilon_0$-moments, $z=\xi^\top S\xi$ and $\|q\|_2^2=\sum_{i=1}^m(\xi_i^2-1)^2=w(\xi)$, so Lemma~\ref{thm:wz_std_lemma} gives
\[
\E[\|q\|_2^2]=2m,\qquad
\E[\|q\|_2^2 z]=(2m+8)\tr(S),
\]
\[
\E[\|q\|_2^2 z^2]
=(2m + 16)\,\tr(S)^2+(4m+32)\,\tr(S^2) + 24 \sum_{i=1}^m S_{ii}^2.
\]
Finally, for $y=s_0^\top M_{se}\varepsilon_0=s_0^\top M_{se}\Sigma^{1/2}\xi$, define
\[
u(s_0):=\Sigma^{1/2}M_{se}^\top s_0\in\mathbb R^m,
\qquad
y=u(s_0)^\top \xi.
\]
Then $y^2=(u(s_0)^\top\xi)^2$ and Lemma~\ref{thm:wz_std_lemma} implies
\begin{align*}
\E[\|q\|_2^2\,y^2]
&=\E_{s_0}\!\left[\E_{\xi}\!\left[w(\xi)\,(\xi^\top u(s_0))^2\mid s_0\right]\right]
=\E_{s_0}\!\left[(2m+8)\,\|u(s_0)\|^2\right]\\
&=(2m+8)\,\E_{s_0}\!\left[s_0^\top(M_{se}\Sigma M_{se}^\top)s_0\right]
=(2m+8)\,\tr\!\big(\Sigma_0\,M_{se}\Sigma M_{se}^\top\big)\\
&=(2m+8)\,\mathrm{IS}_{\Sigma_0}\!\big(M_{se}\Sigma M_{se}^\top\big),
\end{align*}
concluding the proof.
\end{proof}
